\lysh{35296}{3}

\begin{alist}
\item Είναι:
\begin{align*}
|x - 1| \ge 5 \ &\Leftrightarrow \\
\Leftrightarrow (x - 1 \le -5 \text{ ή } x - 1 \ge 5) \ &\Leftrightarrow \\
\Leftrightarrow (x &\le -4 \text{ ή } x \ge 6) \Leftrightarrow x \in (-\infty, -4] \cup [6, +\infty).
\end{align*}

\item Ισχύει ότι:
\begin{align*}
d(x, 5) < 3 \ &\Leftrightarrow 
 |x - 5| < 3 \ \Leftrightarrow \\
\Leftrightarrow -3 < x - 5 < 3 \ &\Leftrightarrow 
 -3 + 5 < x - 5 + 5 < 3 + 5 \ \Leftrightarrow \\
\Leftrightarrow 2 < x < 8 &\Leftrightarrow x \in (2, 8).
\end{align*}

\item Παριστάνουμε τις λύσεις των ανισώσεων στον ίδιο άξονα αριθμών:

\begin{center}
\begin{tikzpicture}
\draw[->] (-5,0) -- (9,0) node[right] {$x$};
\foreach \x/\xtext in {-4/-4, 2/2, 6/6, 8/8}
  \draw[shift={(\x,0)}] (0pt,2pt) -- (0pt,-2pt) node[below] {$\xtext$};

% Λύσεις ανίσωσης α: (-inf, -4] U [6, +inf)
\draw[thick, maincolor] (-5,0.2) -- (-4,0.2);
\filldraw[maincolor] (-4,0.2) circle (2pt);
\draw[thick, maincolor] (6,0.2) -- (9,0.2);
\filldraw[maincolor] (6,0.2) circle (2pt);

% Λύσεις ανίσωσης β: (2, 8)
\draw[thick, secondarycolor] (2,0.4) -- (8,0.4);
\draw[secondarycolor, fill=white] (2,0.4) circle (2pt);
\draw[secondarycolor, fill=white] (8,0.4) circle (2pt);

% Κοινές λύσεις: [6, 8)
\draw[very thick, green!60!black] (6,0.2) -- (8,0.2);
\end{tikzpicture}
\end{center}

Όπως φαίνεται από το σχήμα, οι κοινές λύσεις των δύο ανισώσεων είναι:
\[6 \le x < 8 \Leftrightarrow x \in [6, 8).\]
\end{alist}

